\documentclass[11pt,a4paper]{article}

% ================ Encodage et langue ================
\usepackage[utf8]{inputenc}
\usepackage[T1]{fontenc}
\usepackage[french]{babel}
\usepackage{lmodern}
\usepackage{textcomp}

% ================ Mise en page ================
\usepackage[margin=2.2cm]{geometry}
\usepackage{parskip}
\setlength{\parindent}{0pt}
\setlength{\parskip}{0.6em}

% ================ Couleurs et liens ================
\usepackage[table]{xcolor}
\definecolor{primary}{HTML}{0B5394}
\definecolor{accent}{HTML}{1F77B4}
\definecolor{light}{HTML}{F2F2F2}
\definecolor{codebg}{HTML}{F6F8FA}
\definecolor{rulec}{HTML}{D0D7DE}

\usepackage[colorlinks=true, linkcolor=primary, urlcolor=primary, citecolor=primary]{hyperref}

% ================ Tableaux ================
\usepackage{longtable, booktabs, array, tabularx}
\renewcommand{\arraystretch}{1.25}
\arrayrulecolor{rulec}

% ================ Listes ================
\usepackage{enumitem}
\setlist{itemsep=2pt, topsep=4pt}

% ================ Code et citations ================
\usepackage{fancyvrb}
\usepackage{framed}
\definecolor{shadecolor}{named}{light}

% Encadrés pour les blocs > (quotes)
\usepackage{mdframed}
\newmdenv[
  linecolor=primary,
  linewidth=2pt,
  topline=false, bottomline=false, rightline=false,
  innerleftmargin=10pt, innerrightmargin=8pt,
  innertopmargin=6pt, innerbottommargin=6pt,
  backgroundcolor=codebg,
  skipabove=8pt, skipbelow=8pt
]{quotebox}
\renewenvironment{quote}{\begin{quotebox}}{\end{quotebox}}

% ================ Titres ================
\usepackage{titlesec}
\setcounter{secnumdepth}{0}
\titleformat{\section}{\Large\bfseries\color{primary}}{\thesection.}{0.6em}{}
\titleformat{\subsection}{\large\bfseries\color{accent}}{\thesubsection}{0.6em}{}
\titleformat{\subsubsection}{\normalsize\bfseries}{\thesubsubsection}{0.6em}{}

% ================ En-têtes et pieds de page ================
\usepackage{fancyhdr}
\pagestyle{fancy}
\fancyhf{}
\fancyhead[L]{\small\color{gray}ESPRIT School of Business --- L3 Business Computing BI --- Projet Intégré}
\fancyfoot[C]{\small\thepage}
\renewcommand{\headrulewidth}{0.3pt}
\renewcommand{\footrulewidth}{0pt}

% ================ Pandoc helpers ================
\providecommand{\tightlist}{\setlength{\itemsep}{0pt}\setlength{\parskip}{0pt}}
\providecommand{\passthrough}[1]{#1}

\usepackage{calc}
\providecommand{\real}[1]{#1}


% ================ Métadonnées ================
\title{\color{primary}\textbf{Timeline du Projet --- 5 Jalons}}
\author{Aymen Ben Brik et Hela Hammami}
\date{2026}

\begin{document}

\begin{titlepage}
\centering
\vspace*{1.5cm}

{\Large\bfseries ESPRIT School of Business}\\[0.3cm]
{\large Licence 3 --- Business Computing BI}\\[2cm]

{\color{gray}\large Projet Intégré}\\[0.3cm]
{\color{gray}\large Analyse et Prédiction du Churn Client}\\[2.2cm]

\rule{0.85\linewidth}{1.2pt}\\[0.6cm]
{\Huge\bfseries\color{primary} Timeline du Projet --- 5 Jalons\par}
\rule{0.85\linewidth}{1.2pt}\\[2cm]

{\large Tuteurs du projet}\\[0.2cm]
{\Large\bfseries Aymen Ben Brik et Hela Hammami}\\[1cm]
{\large 2026}

\vfill
{\color{gray}\small Document à destination des étudiants}
\end{titlepage}

\tableofcontents
\newpage

Le projet se déroule du \textbf{9/11 au 7/12/2026} et s'organise en
\textbf{cinq phases}. Chaque phase se termine par un jalon qui
correspond à un point de suivi avec les tuteurs. L'organisation interne
reste libre, mais les échéances sont fermes.

\section{Vue d'ensemble}\label{vue-densemble}

{\def\LTcaptype{none} % do not increment counter
\begin{longtable}[]{@{}
  >{\raggedright\arraybackslash}p{(\linewidth - 6\tabcolsep) * \real{0.2500}}
  >{\raggedright\arraybackslash}p{(\linewidth - 6\tabcolsep) * \real{0.2500}}
  >{\raggedright\arraybackslash}p{(\linewidth - 6\tabcolsep) * \real{0.2500}}
  >{\raggedright\arraybackslash}p{(\linewidth - 6\tabcolsep) * \real{0.2500}}@{}}
\toprule\noalign{}
\begin{minipage}[b]{\linewidth}\raggedright
Phase
\end{minipage} & \begin{minipage}[b]{\linewidth}\raggedright
Échéance
\end{minipage} & \begin{minipage}[b]{\linewidth}\raggedright
Objectif principal
\end{minipage} & \begin{minipage}[b]{\linewidth}\raggedright
Rendu du jalon
\end{minipage} \\
\midrule\noalign{}
\endhead
\bottomrule\noalign{}
\endlastfoot
\textbf{1. Données} & Semaine du 9/11 & Comprendre et qualifier les
données & Notebook d'analyse univariée + décompte des manquants et des
outliers \\
\textbf{2. DWH / ETL} & Semaine du 16/11 & Concevoir et alimenter
l'entrepôt & Modèle en étoile + jobs Talend + base PostgreSQL
alimentée \\
\textbf{3. Reporting} & 23/11 & Restituer les KPI & Rapport Power BI +
mesures DAX documentées \\
\textbf{4. Prédiction} & 1/12 & Prévoir un KPI et prédire le churn & KPI
fixé + modèle LSTM/RNN + modèle de churn évalué \\
\textbf{5. Déploiement} & 7/12 & Livrer l'application & Interface web +
diagrammes UML + rapport et présentation \\
\end{longtable}
}

\begin{center}\rule{0.5\linewidth}{0.5pt}\end{center}

\section{Phase 1 --- Collecte et compréhension des données (semaine du
9/11)}\label{phase-1-collecte-et-compruxe9hension-des-donnuxe9es-semaine-du-911}

\subsection{À réaliser}\label{uxe0-ruxe9aliser}

\begin{itemize}
\tightlist
\item
  Constituer l'équipe, créer le dépôt GitHub sans données, préparer
  l'environnement.
\item
  Charger les fichiers et comprendre le grain, les clés et les fichiers
  de référence.
\item
  Définir le churn et le grain d'analyse.
\item
  Réaliser l'\textbf{analyse univariée}.
\item
  Produire le \textbf{décompte des valeurs manquantes} et des
  \textbf{outliers}.
\end{itemize}

\subsection{Bonus}\label{bonus}

\begin{itemize}
\tightlist
\item
  \textbf{Web scraping} avec Beautiful Soup, Playwright ou Scrapy.
\item
  \textbf{Analyse bivariée}.
\end{itemize}

\subsection{Rendu}\label{rendu}

\begin{itemize}
\tightlist
\item
  Notebook documenté et tableau de synthèse par variable : type, \% de
  manquants, nombre d'outliers, méthode, décision proposée.
\end{itemize}

\begin{center}\rule{0.5\linewidth}{0.5pt}\end{center}

\section{Phase 2 --- Conception du DWH et ETL (semaine du
16/11)}\label{phase-2-conception-du-dwh-et-etl-semaine-du-1611}

\subsection{À réaliser}\label{uxe0-ruxe9aliser-1}

\begin{itemize}
\tightlist
\item
  Concevoir le modèle en étoile : faits, dimensions, grain, dimension
  temps.
\item
  Créer le schéma sous \textbf{PostgreSQL} avec \textbf{pgAdmin}.
\item
  Développer les jobs \textbf{Talend} : extraction, nettoyage,
  jointures, chargement.
\item
  Contrôler la cohérence entre la source et l'entrepôt.
\end{itemize}

\subsection{Bonus}\label{bonus-1}

\begin{itemize}
\tightlist
\item
  \textbf{SSIS} ; \textbf{Databricks} en architecture bronze / silver /
  gold ; \textbf{SCD de type 2} ; \textbf{Python}.
\end{itemize}

\subsection{Rendu}\label{rendu-1}

\begin{itemize}
\tightlist
\item
  Diagramme en étoile, scripts SQL, jobs Talend exportés, requêtes de
  contrôle.
\end{itemize}

\begin{center}\rule{0.5\linewidth}{0.5pt}\end{center}

\section{Phase 3 --- Reporting et dashboard (pour le
23/11)}\label{phase-3-reporting-et-dashboard-pour-le-2311}

\subsection{À réaliser}\label{uxe0-ruxe9aliser-2}

\begin{itemize}
\tightlist
\item
  Connecter \textbf{Power BI} à l'entrepôt.
\item
  Écrire les KPI en \textbf{DAX}.
\item
  Construire au moins trois pages : vue d'ensemble, analyse du churn,
  segmentation client.
\end{itemize}

\subsection{Bonus}\label{bonus-2}

\begin{itemize}
\tightlist
\item
  Visualisation avec \textbf{Databricks} ; \textbf{interface
  JavaScript}.
\end{itemize}

\subsection{Rendu}\label{rendu-2}

\begin{itemize}
\tightlist
\item
  Fichier \texttt{.pbix} remis aux tuteurs hors GitHub, captures sans
  ligne client et liste des mesures DAX documentées dans le dépôt.
\end{itemize}

\begin{center}\rule{0.5\linewidth}{0.5pt}\end{center}

\section{Phase 4 --- Modèles prédictifs (pour le
1/12)}\label{phase-4-moduxe8les-pruxe9dictifs-pour-le-112}

\subsection{À réaliser}\label{uxe0-ruxe9aliser-3}

\begin{itemize}
\tightlist
\item
  \textbf{Fixer un KPI} et construire sa série temporelle.
\item
  \textbf{Analyser la série} : tendance, saisonnalité, pics,
  stationnarité.
\item
  Entraîner un modèle \textbf{LSTM ou RNN} et le comparer à une
  référence simple.
\item
  Entraîner et comparer des modèles de \textbf{Machine Learning} pour le
  churn, après contrôle de la fuite de cible.
\end{itemize}

\subsection{Bonus}\label{bonus-3}

\begin{itemize}
\tightlist
\item
  Appliquer la démarche à \textbf{plusieurs KPI}.
\end{itemize}

\subsection{Rendu}\label{rendu-3}

\begin{itemize}
\tightlist
\item
  Notebooks complets, modèles sauvegardés hors du dépôt et tableaux
  comparatifs.
\end{itemize}

\begin{center}\rule{0.5\linewidth}{0.5pt}\end{center}

\section{Phase 5 --- Déploiement (pour le
7/12)}\label{phase-5-duxe9ploiement-pour-le-712}

\subsection{À réaliser}\label{uxe0-ruxe9aliser-4}

\begin{itemize}
\tightlist
\item
  Développer et déployer l'\textbf{interface web}.
\item
  Produire les \textbf{diagrammes UML} : cas d'utilisation global, une
  séquence, classes.
\item
  Finaliser le rapport, les slides et le dépôt GitHub.
\end{itemize}

\subsection{Bonus}\label{bonus-4}

\begin{itemize}
\tightlist
\item
  \textbf{DevOps} : tests automatisés ; \textbf{application mobile}.
\end{itemize}

\subsection{Rendu}\label{rendu-4}

\begin{itemize}
\tightlist
\item
  URL de l'application, dépôt finalisé, rapport PDF (25--40 pages) et
  slides.
\end{itemize}

\begin{center}\rule{0.5\linewidth}{0.5pt}\end{center}

\section{Recommandations transverses}\label{recommandations-transverses}

\begin{itemize}
\tightlist
\item
  \textbf{Traitez le socle d'abord} : un bonus ne remplace pas une
  partie obligatoire.
\item
  \textbf{Commits réguliers} : chaque membre commit sur sa branche tout
  au long du projet.
\item
  \textbf{Branches Git} : utilisez \texttt{main} pour le code stable et
  des branches \texttt{feature/...} par tâche.
\item
  \textbf{Documentation au fil de l'eau} : rédigez le rapport à chaque
  jalon, pas à la fin.
\item
  \textbf{Marge de déploiement} : la phase 4 et la phase 5 sont
  rapprochées. Testez un premier déploiement avant le 1/12.
\end{itemize}

\end{document}
